\item La intensidad $I$ de la luz que llega al CCD en una cámara es proporcional al área de la lente y al cuadrado de su diámetro $D$, e inversamente proporcional al cuadrado de la distancia focal $f$. La relación $f/D$ se denomina el número $f$ de la lente. Suponga que la lente de una cámara digital tiene una distancia focal de $55\text{ mm}$ y una velocidad de $f/1.8$. Bajo ciertas condiciones, el tiempo de exposición óptimo es de $1/500\text{ s}$.
\begin{enumerate}
    \item Determine el diámetro de apertura de la lente.
    \item Calcule el tiempo de exposición óptimo si el número $f$ se cambia a $f/4.0$ bajo las mismas condiciones de iluminación.
\end{enumerate}